% ============================================================
\chapter{Courbures Sectionnelle, de Ricci, Scalaire}
\label{ch:courbures}
% ============================================================

Le tenseur de Riemann est un objet riche mais complexe --- en dimension $n$, il poss\`ede $\frac{n^2(n^2-1)}{12}$ composantes ind\'ependantes. Pour extraire une information g\'eom\'etrique plus lisible, on contracte ce tenseur de diff\'erentes mani\`eres. La \emph{courbure sectionnelle} mesure la courbure dans chaque plan tangent --- c'est la g\'en\'eralisation directe de la courbure de Gauss des surfaces. La \emph{courbure de Ricci}, obtenue par moyenne sur les plans contenant une direction donn\'ee, contr\^ole la convergence ou la divergence des g\'eod\'esiques et appara\^it dans les \'equations d'Einstein. La \emph{courbure scalaire}, derni\`ere contraction, donne un unique nombre en chaque point. Ces trois niveaux de lecture forment la hi\'erarchie fondamentale de la g\'eom\'etrie riemannienne.

\section{Courbure sectionnelle}

\begin{definition}[Courbure sectionnelle]
Soit $(M, g)$ une vari\'et\'e riemannienne et $\sigma \subset T_pM$ un $2$-plan
engendr\'e par deux vecteurs lin\'eairement ind\'ependants $u, v$. La \emph{courbure
sectionnelle} de $\sigma$ est~:
\[
    K(\sigma) = K(u, v) = \frac{\inner{R(u,v)v}{u}}{\norm{u}^2\norm{v}^2 - \inner{u}{v}^2}.
\]
\end{definition}

\begin{proposition}
La courbure sectionnelle $K(u,v)$ ne d\'epend que du $2$-plan $\sigma = \operatorname{Vect}(u,v)$
et non du choix de base.
\end{proposition}

\begin{proof}
Si $(u', v') = (au + bv, cu + dv)$ est une autre base de $\sigma$ avec $ad - bc \neq 0$,
alors par les sym\'etries de $R$~:
$\inner{R(u',v')v'}{u'} = (ad-bc)^2 \inner{R(u,v)v}{u}$
et $\norm{u'}^2\norm{v'}^2 - \inner{u'}{v'}^2 = (ad-bc)^2(\norm{u}^2\norm{v}^2 - \inner{u}{v}^2)$.
Le quotient est donc invariant.
\end{proof}

\begin{theoreme}[D\'etermination du tenseur de Riemann]
La courbure sectionnelle d\'etermine compl\`etement le tenseur de Riemann.
Plus pr\'ecis\'ement, si $K(u,v)$ est connu pour tous les $2$-plans, alors
$R$ est uniquement d\'etermin\'e.
\end{theoreme}

\begin{proof}
La fonction $(u,v) \mapsto \inner{R(u,v)v}{u}$ est une forme quadratique en
chacune des variables $u$ et $v$. Par les sym\'etries de $R$, la connaissance
de cette forme pour tous $u, v$ d\'etermine $R_{ijkl}$ par polarisation.
\end{proof}

\section{Espaces \`a courbure constante}

\begin{definition}
$(M, g)$ est un \emph{espace \`a courbure (sectionnelle) constante} $\kappa$ si
$K(\sigma) = \kappa$ pour tout $2$-plan $\sigma$ en tout point.
\end{definition}

\begin{theoreme}[Forme du tenseur de Riemann]
Si $K \equiv \kappa$, alors~:
\[
    R(X,Y)Z = \kappa\big(\inner{Y}{Z}X - \inner{X}{Z}Y\big).
\]
En coordonn\'ees~: $R_{ijkl} = \kappa(g_{ik}g_{jl} - g_{il}g_{jk})$.
\end{theoreme}

\begin{formules}[title=\textbf{Espaces \`a courbure constante}]
\begin{center}
\renewcommand{\arraystretch}{1.3}
\begin{tabular}{lccc}
\hline
\textbf{Espace} & $\kappa$ & \textbf{Mod\`ele} & \textbf{Topologie} \\
\hline
Sph\`ere $S^n$ & $+1$ & $\{x \in \R^{n+1} : \norm{x} = 1\}$ & compacte, $\pi_1 = 0$ \\
Euclidien $\R^n$ & $0$ & $\R^n$ & non compacte, $\pi_1 = 0$ \\
Hyperbolique $\mathbb{H}^n$ & $-1$ & demi-espace & non compacte, $\pi_1 = 0$ \\
\hline
\end{tabular}
\end{center}
\end{formules}

\begin{theoreme}[Classification -- Killing--Hopf]
Soit $(M, g)$ une vari\'et\'e riemannienne compl\`ete simplement connexe de courbure
sectionnelle constante $\kappa$. Alors $(M, g)$ est isom\'etrique \`a~:
\begin{itemize}
    \item $S^n(1/\sqrt{\kappa})$ si $\kappa > 0$,
    \item $\R^n$ si $\kappa = 0$,
    \item $\mathbb{H}^n(1/\sqrt{-\kappa})$ si $\kappa < 0$.
\end{itemize}
Toute vari\'et\'e compl\`ete de courbure constante est un quotient de ces espaces par
un sous-groupe discret d'isom\'etries agissant librement et proprement.
\end{theoreme}

\section{Courbure de Ricci}

\begin{definition}[Tenseur de Ricci]
Le \emph{tenseur de Ricci} est la contraction du tenseur de Riemann~:
\[
    \operatorname{Ric}(X, Y) = \tr\big(Z \mapsto R(Z, X)Y\big).
\]
En coordonn\'ees~: $R_{ij} = R^k_{ikj} = g^{kl} R_{kilj}$.
\end{definition}

\begin{proposition}
Le tenseur de Ricci est sym\'etrique~: $\operatorname{Ric}(X,Y) = \operatorname{Ric}(Y,X)$,
i.e.\ $R_{ij} = R_{ji}$.
\end{proposition}

\begin{proof}
$R_{ij} = g^{kl}R_{kilj}$. Par la sym\'etrie par paires $R_{kilj} = R_{ljki}$,
et par l'antisym\'etrie $R_{ljki} = R_{jlki} \cdot (-1) \cdot (-1) = R_{ljki}$.
On obtient $R_{ij} = R_{ji}$ par la sym\'etrie $R_{kilj} = R_{kjli}$... plus
directement, en base orthonorm\'ee~:
$R_{ij} = \sum_k R_{kikj}$ et $R_{ji} = \sum_k R_{kjki} = \sum_k R_{kikj} = R_{ij}$.
\end{proof}

\begin{proposition}[Interpr\'etation g\'eom\'etrique]
Si $(e_1, \ldots, e_n)$ est une base orthonorm\'ee de $T_pM$ avec $e_1 = v/\norm{v}$,
alors~:
\[
    \operatorname{Ric}(v, v) = \norm{v}^2 \sum_{i=2}^n K(v, e_i).
\]
La courbure de Ricci dans la direction $v$ est la \emph{somme des courbures
sectionnelles} des $2$-plans contenant $v$.
\end{proposition}

\begin{exemple}[Ricci de $S^n$]
Toutes les courbures sectionnelles valent $1$, donc~:
$\operatorname{Ric}(v,v) = (n-1)\norm{v}^2$, i.e.\ $\operatorname{Ric} = (n-1)g$.
\end{exemple}

\begin{exemple}[Ricci de $\mathbb{H}^n$]
$\operatorname{Ric} = -(n-1)g$.
\end{exemple}

\begin{exemple}[Ricci d'un groupe de Lie bi-invariant]
Avec $R(X,Y)Z = \frac{1}{4}[[X,Y],Z]$, on obtient~:
\[
    \operatorname{Ric}(X,Y) = -\frac{1}{4}\tr(\operatorname{ad}_X \circ \operatorname{ad}_Y)
    = -\frac{1}{4}B(X,Y),
\]
o\`u $B$ est la forme de Killing de $\mathfrak{g}$.
\end{exemple}

\section{Courbure scalaire}

\begin{definition}[Courbure scalaire]
La \emph{courbure scalaire} est la trace du tenseur de Ricci~:
\[
    S = \operatorname{Scal} = \tr_g \operatorname{Ric} = g^{ij} R_{ij}.
\]
Si $(e_1, \ldots, e_n)$ est une base orthonorm\'ee~:
$S = \sum_{i=1}^n \operatorname{Ric}(e_i, e_i) = \sum_{i \neq j} K(e_i, e_j)
= 2\sum_{i < j} K(e_i, e_j)$.
\end{definition}

\begin{exemple}
Pour $S^n$~: $S = n(n-1)$. Pour $\mathbb{H}^n$~: $S = -n(n-1)$. Pour $\R^n$~: $S = 0$.
\end{exemple}

\section{Vari\'et\'es d'Einstein}

\begin{definition}[Vari\'et\'e d'Einstein]
$(M, g)$ est une \emph{vari\'et\'e d'Einstein} si le tenseur de Ricci est
proportionnel \`a la m\'etrique~:
\[
    \operatorname{Ric} = \lambda\, g
\]
pour une constante $\lambda \in \R$. En prenant la trace~: $S = n\lambda$,
d'o\`u $\lambda = S/n$.
\end{definition}

\begin{exemple}
Les espaces \`a courbure constante sont des vari\'et\'es d'Einstein avec
$\lambda = (n-1)\kappa$.
\end{exemple}

\begin{theoreme}[Schur]
Si $n \geq 3$ et la courbure sectionnelle est constante en chaque point
(i.e.\ $K_p(\sigma) = f(p)$ pour tout $2$-plan $\sigma$ en $p$), alors $f$ est
constante sur $M$ (i.e.\ $(M,g)$ est \`a courbure constante).

Plus g\'en\'eralement, si $n \geq 3$ et $\operatorname{Ric} = f\, g$ avec
$f \in C^\infty(M)$, alors $f$ est constante.
\end{theoreme}

\begin{proof}
De $\operatorname{Ric} = f\, g$, on tire $S = nf$ et la seconde identit\'e de
Bianchi contract\'ee donne~:
\[
    \nabla^j R_{ij} = \frac{1}{2}\nabla_i S.
\]
Or $\nabla^j R_{ij} = \nabla^j(fg_{ij}) = (\nabla_j f) g^{ij}\cdot g_{ij} = \nabla_i f$.
D'o\`u $\nabla_i f = \frac{1}{2}\nabla_i(nf) = \frac{n}{2}\nabla_i f$, soit
$(1 - n/2)\nabla_i f = 0$. Pour $n \geq 3$, cela implique $\nabla f = 0$.
\end{proof}

\section{Tenseur de Weyl}

\begin{definition}[Tenseur de Weyl]
Pour $n \geq 3$, le \emph{tenseur de Weyl} $W$ est la partie du tenseur de Riemann
qui n'est pas d\'etermin\'ee par Ricci~:
\begin{align*}
    R_{ijkl} &= W_{ijkl} + \frac{1}{n-2}\left(R_{ik}g_{jl} - R_{il}g_{jk}
    + R_{jl}g_{ik} - R_{jk}g_{il}\right) \\
    &\quad - \frac{S}{(n-1)(n-2)}\left(g_{ik}g_{jl} - g_{il}g_{jk}\right).
\end{align*}
\end{definition}

\begin{proposition}
Le tenseur de Weyl est invariant conforme~: si $\tilde{g} = e^{2\varphi}g$, alors
$\tilde{W}^i_{jkl} = W^i_{jkl}$.
\end{proposition}

\begin{proposition}
$W \equiv 0$ si et seulement si $n \leq 3$, ou si $n \geq 4$ et $(M,g)$ est
\emph{conform\'ement plate} (localement conforme \`a $\R^n$).
\end{proposition}

\section{Calculs sur les produits tordus}

\begin{proposition}[Courbure du produit tordu]
Sur $B \times_f F$ avec $g = g_B + f^2 g_F$, si $\dim B = 1$ (i.e.\ $B = I \subset \R$,
$g_B = dt^2$)~:
\begin{enumerate}
    \item Pour $V, W$ tangents \`a $F$~:
    $K(V, W) = \frac{K_F(V,W) - (f')^2}{f^2}$.
    \item Pour $V$ tangent \`a $F$ et $\partial_t$ tangent \`a $B$~:
    $K(\partial_t, V) = -\frac{f''}{f}$.
\end{enumerate}
\end{proposition}

\begin{exemple}[Courbure de $\R^n$ en coordonn\'ees polaires]
$g = dr^2 + r^2 g_{S^{n-1}}$, donc $f(r) = r$, $f' = 1$, $f'' = 0$.
$K(\partial_r, V) = 0$ et $K(V, W) = \frac{1 - 1}{r^2} = 0$. Tout est plat.
\end{exemple}

\begin{exemple}[Courbure de $S^n$ en coordonn\'ees g\'eod\'esiques]
$g = dr^2 + \sin^2\!r\, g_{S^{n-1}}$, $f(r) = \sin r$, $f'' = -\sin r$.
$K(\partial_r, V) = -\frac{-\sin r}{\sin r} = 1$ et
$K(V,W) = \frac{1 - \cos^2 r}{\sin^2 r} = 1$.
\end{exemple}

\begin{exemple}[Courbure de $\mathbb{H}^n$ en coordonn\'ees g\'eod\'esiques]
$g = dr^2 + \sinh^2\!r\, g_{S^{n-1}}$, $f(r) = \sinh r$, $f'' = \sinh r$.
$K(\partial_r, V) = -\frac{\sinh r}{\sinh r} = -1$ et
$K(V,W) = \frac{1 - \cosh^2 r}{\sinh^2 r} = -1$.
\end{exemple}

\section{Identit\'e de Bianchi contract\'ee}

\begin{theoreme}[Identit\'e de Bianchi contract\'ee]
\[
    \nabla^j R_{ij} = \frac{1}{2}\nabla_i S,
\]
ou de mani\`ere \'equivalente, le \emph{tenseur d'Einstein}
$G_{ij} = R_{ij} - \frac{1}{2}Sg_{ij}$ est \`a divergence nulle~:
$\nabla^j G_{ij} = 0$.
\end{theoreme}

\begin{remarque}
Cette identit\'e est fondamentale en relativit\'e g\'en\'erale, o\`u les \'equations
d'Einstein s'\'ecrivent $G_{ij} = 8\pi T_{ij}$, et la conservation de
l'\'energie-impulsion $\nabla^j T_{ij} = 0$ est automatique.
\end{remarque}

\section{Exercices}

\begin{exercice}
Calculer la courbure sectionnelle, le tenseur de Ricci et la courbure scalaire
du tore $S^1(r_1) \times S^1(r_2)$ avec la m\'etrique produit.
\end{exercice}

\begin{exercice}
Montrer que si $(M, g)$ est de dimension $2$, alors
$\operatorname{Ric} = K\, g$ et $S = 2K$.
\end{exercice}

\begin{exercice}
Montrer que $\C P^n$ muni de la m\'etrique de Fubini--Study est une vari\'et\'e
d'Einstein avec $\lambda = 2(n+1)$.
\end{exercice}

\begin{exercice}
Calculer le tenseur de Weyl de $S^2 \times S^2$ et montrer qu'il n'est pas
conform\'ement plat pour $n = 4$.
\end{exercice}

\begin{exercice}
Pour le produit tordu $g = dt^2 + e^{2t} g_{\R^{n-1}}$, calculer toutes les
courbures sectionnelles, le tenseur de Ricci et la courbure scalaire.
\end{exercice}

\begin{exercice}
Montrer le th\'eor\`eme de Schur en dimension $3$ en utilisant la seconde identit\'e
de Bianchi. Donner un contre-exemple montrant que le th\'eor\`eme est faux en
dimension $2$.
\end{exercice}
